% ------------------------------------------------------------------
% Sección 5.9: Conclusiones del Capítulo 5
% ------------------------------------------------------------------
\section{Conclusiones del Capítulo}
\label{sec:cap5_conclusiones}

Este capítulo ha transitado del dominio de la predicción técnica al de la acción operativa, cerrando el ciclo de diseño iniciado en el Capítulo~\ref{chap:cap2} y consolidado en el Capítulo~\ref{chap:cap4}. Si el Capítulo~\ref{chap:cap4} demostró que un sistema híbrido DEB--MLP puede estimar emisiones de CO$_2$ con MAPE $= 8{,}3$\% y detectar eventos de anaerobiosis con Recall $= 100$\%, el presente capítulo ha establecido que esas predicciones adquieren valor ambiental solo cuando se traducen en decisiones de manejo concretas, trazables y verificables.

\textbf{Validación predictiva.} El sistema híbrido superó todos los umbrales de éxito establecidos en el proyecto: $R^2 = 0{,}78 \geq 0{,}70$ (H2), MAPE $= 8{,}3$\% $< 10$\% (H3), y Recall $= 100$\% $> 80$\% en detección de anaerobiosis (H4). La discrepancia acumulada de carbono $\bar{\Delta}_{\mathrm{Total}} = +5{,}1$~g~C (3{,}6\% del total medido) se reportó explícitamente como métrica de incertidumbre modelada, no como error a ocultar. El análisis de residuos confirmó ausencia de estructura temporal no capturada y distribución compatible con normalidad, validando la especificación del modelo.

\textbf{Plataforma de soporte a la decisión.} La arquitectura de tres capas propuesta ---percepción con edge computing, procesamiento híbrido y aplicación con dashboard web--- responde a las restricciones contextuales del sur global: conectividad intermitente, costos marginales ajustados y soberanía de datos. El nodo sensor basado en ESP32 con un único sensor NDIR de CO$_2$ (rango 6000~ppm) y sensores complementarios de CH$_4$, temperatura y humedad, demuestra que la precisión instrumental no requiere hardware de laboratorio, sino protocolos de calibración rigurosos y compensación cruzada de señal. La autonomía de 72~horas del buffer local y la ejecución de algoritmos de detección en el borde garantizan funcionalidad continua incluso en zonas rurales sin cobertura estable.

\textbf{Semáforo de riesgo ambiental.} El mecanismo central de traducción entre modelo y operador clasifica el estado del reactor en tres niveles (verde, ámbar, rojo) mediante una lógica multivariable que combina la discrepancia CO$_2$ medido versus estimado, la tendencia de CH$_4$ y la pendiente de la relación CO$_2$/O$_2$. Esta estrategia evita la simplificación peligrosa de umbrales absolutos y respeta la complejidad biológica del proceso. Las acciones asociadas a cada nivel (continuidad, revisión de humedad, volteo inmediato) son deliberadamente de bajo costo mecánico, ejecutables sin herramientas especializadas ni energía externa, lo cual asegura su viabilidad para pequeños productores.

\textbf{Protocolo de respuesta ante anaerobiosis.} Se formalizó una secuencia de cuatro pasos con tiempos cuantificados: confirmación (0--2~min), preparación (2--5~min), volteo (5--15~min) y verificación (15--30~min). La tasa de mitigación $\eta_{\mathrm{mit}} > 80$\%, validada experimentalmente en cuatro eventos confirmados, demuestra que la intervención física dirigida por el sistema reduce efectivamente las emisiones fugitivas de CH$_4$. La inclusión de una regla de escalamiento a supervisor técnico cuando la humedad supera el 80\% reconoce los límites de la autonomía operativa y establece un mecanismo de gobernanza ambiental escalonada.

\textbf{Generación automatizada de reportes MRV.} El sistema se alinea con estándares internacionales (IPCC Tier 3, Verra VCS v4.7) mediante una estructura de cinco bloques: inventario dinámico de carbono con discrepancia explícita ($\Delta_{\mathrm{Total}} = +5{,}1$~g~C), métricas predictivas del modelo, log de eventos de anaerobiosis, línea base de referencia y metadatos completos del ciclo. Los mecanismos de timestamping estricto, hash encadenado por bloques y registro de provenance computacional garantizan la auditabilidad y la detección de alteraciones, respondiendo al riesgo de \emph{greenwashing} digital identificado en el marco teórico.

\textbf{Utilidad operativa.} Las métricas de usabilidad para no expertos, latencia end-to-end $< 1$~min en conectividad estable y pertinencia contextual con costo marginal $< 5$\% validan que la plataforma no es un artefacto de laboratorio, sino una herramienta de gestión ambiental integrable en prácticas productivas reales. El cumplimiento de estos umbrales habilita la transición del prototipo (TRL 5) hacia una tecnología validada en entorno operativo (TRL 6).

En síntesis, este capítulo ha transformado el sistema híbrido DEB--MLP de un instrumento de medición en una infraestructura de soporte a la decisión. La predicción matemática, la detección de anomalías y la cuantificación de carbono se materializan ahora en acciones de manejo, registros auditables y reportes comparables con estándares globales. Esta materialización no sustituye el juicio técnico del productor, pero lo estructura, externaliza y verifica, fortaleciendo la autonomía técnica del sur global frente a los mercados de carbono y las políticas de economía circular.

El Capítulo~\ref{chap:cap6}, que cierra esta tesis, sintetizará los aportes integrales del sistema híbrido, discutirá las limitaciones encontradas y trazará las líneas de trabajo futuro que permitan escalar la bioconversión con \emph{Hermetia illucens} como una tecnología climáticamente verificable y socialmente justa.
